\subsection{Conditional break-even screens}\begin{table}[htbp]\centering\small\begin{tabular}{lrrrr}
\toprule
Condition & Useful / month & Required purchase & Validation ceiling* & Capacity upper*\\\midrule
Conservative & 0.229 & 38.0\% & \usd{3,068} & 13.4\\
Central & 0.086 & 2.8\% & \usd{9,196} & 22.0\\
Favorable & 0.039 & 0.6\% & \usd{22,016} & 29.1\\
\bottomrule\end{tabular}\caption{Analytic screens at the confirmation mean posted price. *Validation ceiling is per useful outcome at one useful project/month; capacity is purchased project attempts/month before bursts, queues, retention and maintenance. Complete setup, settlement-runoff overhead and qualification are included; all false passes are refunded in this screen. These are conditional thresholds, not demand evidence or service guarantees.}\end{table}
The primary retention options cost \usd{290.05} for R, \usd{307.95} for V, \usd{569.72} for bundle. Their optimistic reconstruction-only reuse requirements are 5.2 for R, 8.3 for V, 6.2 for bundle. The source job is not charged again. Actual repayment requires enough compatible first encounters with recipients after overlap, transfer, access, upkeep, erosion and the capacity consumed by retention.

For an additional retention bill $K$, maximum affordable expenditure at a decision is cleared cash minus protected liabilities, reserves and headroom. The full ledger determines this ceiling before spending. For constant compatible reuse $u$ and per-use incremental value $b_t$, the prospective screen is $\sum_{t=1}^{H}\delta^t(ub_t-m_t)>K$, where $m_t$ includes remaining maintenance and revalidation. Advancing Beta changes $b_t$; a shorter useful lifetime or lower adoption can make waiting optimal even after a profitable source job.
